\subsection{无上同调性}\label{sec:17h}
将\bcite{131}引理 2.32 证明中的代数论证与该文第 3 章中的几何论证结合起来，即可证明\eref{17.16}所定义的复形 $D$ 无上同调；鉴于\stref{引理}{5.3}(iii)，这一点对我们很重要。\znote{256}{“无上同调”与“$t$ 是同构”实为一事，桥梁是映射锥：引理 5.3 的证明指出 $t^1$ 的映射锥是复形 $\widetilde D$，其第三个分量把 (5.8) 中的 $d(c_1)+l(c_0,b)$ 换成 $\mu^2_{\A}(c,c_1)+\mu^2_{\A}(k(c_0),b)+\mu^3_{\A}(c,c_0,b)$，而 $(\xi,\psi)$ 正是补出这一替换的链同构。17g 节的 $\tilde d=\mu^2(c,\cdot)$、$\tilde l=\mu^2(k(\cdot),\cdot)+\mu^3(c,\cdot,\cdot)$ 就是这份锥数据，故\eref{17.16}无上同调即 $t$ 为拟同构。注意 $D$ 与 $\widetilde D$ 只有第三个分量不同，其余两项完全一样。[来源：Seidel, Fukaya categories and Picard–Lefschetz theory, §5c 引理 5.3(iii) 及其证明。]}严格说来，\bcite{131}只在若干附加技术条件下证明了无上同调性，即：
\begin{enumerate}[label=(\roman*)]
\item $Q,L_0,L_2$ 处于一般位置。对子流形对 $(Q,L_0)$、$(L_0,L_2)$ 和 $(Q,L_2)$，Floer 数据的 Hamilton 项应为零。此外，对作用量泛函\eref{8.11}在交点处的取值还有一个条件：$x_0\in\C(Q,L_2)$ 时的各数 $A(x_0)$，以及 $(x_1,x_2)\in\C(Q,L_0)\times\C(L_0,L_2)$ 时的各数 $A(x_1)+A(x_2)$，合在一起必须两两不同。
\item 必须能将标架 $L_0=V\cong S^n$ 延拓为余切球丛的辛嵌入 $B_s^*S^n\to\operatorname{int}(M)$，其中某个 $s>0$，并使 $\theta_M$ 的拉回为 $T^*S^n$ 上的标准一形式。此外，在这个管状邻域内，$Q$ 与 $L_2$ 都应为若干余切纤维的并。
\item $L_1=\tau_V^{-1}(L_2)$，其中 Dehn 扭转使用先前选定的管状邻域定义。决定扭转速度的函数 $b$，见\eref{16.5}，必须满足某些条件；这些条件依赖于 (i) 中作用量值之间的最小差，以及点 $x_1,x_2$ 关于 $V$ 上球面度量的最小距离。
\item 对子流形对 $(L_0,L_1)$ 和 $(Q,L_1)$，Floer 数据的 Hamilton 项应为零；类似地，三元组 $(Q,L_0,L_1)$ 的扰动数据应满足 $K=0$。
\item 最后，定义 $(d,l)$ 时使用的辅助几何数据，即 Lefschetz 纤维化、Lagrangian 边界条件及相对扰动数据，必须按一种精确受控的方式选取。
\end{enumerate}
首先观察到，只要允许作下列改变，就能满足所有这些条件：对各 $L_k$ 作 Lagrangian 同痕；将 $\theta_M$ 换为 $\theta_M+dh$，其中函数 $h$ 支撑于 $\operatorname{int}(M)$；以及改变用来定义 Fukaya 范畴的 Floer 数据和扰动数据。这些改变都不影响上同调层次上对象的同构类，无论是在 Fukaya 范畴中，还是在其 $\Ai$ 模范畴中。因此，应用\stref{引理}{5.3}(iii)，已经可以得出 $T_{Y_0}(\Y_1)\cong\Y_2$。对大多数用途而言，这已经足够，但仍不完全令人满意。因此，我们将利用\secref{17f}的结果作更细致的分析。首先，(v) 显然并非必要，因为正如\eref{17.17}所示，$D$ 的同构类与这些选择无关。
\SourcePage{265}
类似地，\eref{17.18}说明它与 $L_1$ 无关，因而可以去掉 (iii)；由于 (ii) 中的管状邻域只用于定义 $L_1=\tau_V^{-1}(L_2)$，这个条件也变得多余。此外，换成与之同痕的 $\widetilde L_1$，也就使得与该子流形相关的 Floer 数据可以完全自由选择，而此前这些数据受 (iv) 限制。对 $L_1$ 所用的论证同样表明，可以在 $L_2$ 的同痕类中移动它，而不影响 $D$ 的拟同构类。最后，正如\stref{引理}{5.3}(iii)中已经提到的，$(d,l)$ 通过\eref{17.19}与 $(c,k)$ 联系这一事实，意味着\eref{17.16}对于某个 $Q$ 无上同调，就对所有与之同痕的 $\widetilde Q$ 也无上同调。合起来，这就使 (i) 以及 (iv) 的剩余部分都不再必要。因此，我们已经证明：\emph{对于任意 $(L_0,L_1,L_2,Q)$，无论定义 $(d,l)$ 时如何作出选择，\eref{17.16}中的复形 $D$ 都无上同调。}用\stref{引理}{5.3}的语言来说，这意味着由 $(c,k)$ 定义的态射 $T_{Y_0}(\Y_1)\to\Y_2$ 始终为同构。

\StatementNumber{13}\begin{note}\label{st:17.13}
作为这一讨论的补充，我们粗略说明怎样构造满足\stref{引理}{5.4}各条件的附加代数数据 $(g,p,q)$。在上同调层次上，$[g]$ 是延拓同构
\[
HF^{\mathrm{pr}}(L_0,L_2)\longrightarrow HF^{\mathrm{pr}}(L_0,\tau_V(L_1))\cong HF^{\mathrm{pr}}(L_0,L_1)
\]
的 Poincaré 对偶。可以用一个技巧简化任务：适当选择 $(L_2,L_0)$ 的 Floer 数据，使其上链群自然地对偶于 $(L_0,L_1)$ 的上链群。若 $L_2$ 就等于 $\tau_V(L_1)$，实现方法是显然的：取 $(L_0,L_1)$ 的 Floer 数据 $(H,J)$，将其换为 $(-H(-t),J(-t))$，再对其施以 $\tau_V$。一般而言，$L_2$ 只与 $\tau_V(L_1)$ 同痕，但可以将这一同痕纳入上述 Floer 数据的修正版中。现在，取 $g$ 为恒等映射 $CF^{\mathrm{pr}}(L_2,L_0)^\vee\to CF^{\mathrm{pr}}(L_0,L_1)$ 的对偶。由此，其余各项及其所需性质可以写成
\begin{equation}\label{eq:17.21}
\boxed{\begin{aligned}
p &:CF^{\mathrm{pr}}(L_2,L_0)\longrightarrow CF^{\mathrm{pr}}(L_1,L_0),\\
q &:CF^{\mathrm{pr}}(L_2,L_0)\otimes CF^{\mathrm{pr}}(L_0,L_1)\longrightarrow CF^{\mathrm{pr}}(L_0,L_0),\\
0&=\mu^1(p(\cdot))+p(\mu^1(\cdot))+\mu^2(\cdot,c),\\
0&=\mu^1(q(\cdot,\cdot))+q(\mu^1(\cdot),\cdot)+q(\cdot,\mu^1(\cdot))\\
 &\quad+\mu^3(\cdot,c,\cdot)+\mu^2(\cdot,k(\cdot))+\mu^2(p(\cdot),\cdot)\\
 &\quad+e\langle\cdot,\cdot\rangle.
\end{aligned}}
\end{equation}
在最后一行中，$\langle\cdot,\cdot\rangle$ 为对偶配对，而 $e\in CF^{\mathrm{pr}}(L_0,L_0)$ 是上同调单位的一个代表。从这些等式的第一个关系式，就能清楚看出如何调整前面 $k$ 的定义来构造 $p$。对于 $q$，考虑\sr{fig:17.9}{图~17.9}中那样的、以 $\RG^{3,1}$ 为底空间的双参数族，不过此时令 $Q=L_0$。边 (i)、(ii) 和 (vi) 的贡献给出\eref{17.21}倒数第二行中的各项，而 (iii) 与 (v) 的贡献消失。于是，
\SourcePage{266}
最后剩下且最有意思的是 (iv)。有关几何设置由两个 Lefschetz 纤维化组成，并将它们在边界标点上方的纤维作一等同，见\sr{fig:17.10}{图~17.10}。左边那一块记为 $E^e\to S^e$，其标点位置可变，因此得到一个带附加参数的求值映射。假设 $(L_0,L_0)$ 的 Floer 数据满足：Hamilton 函数为 $L_0=S^n$ 上的标准 Morse 函数，从而 $CF^{\mathrm{pr}}(L_0,L_0)$ 有两个生成元 $e=y_0$ 与 $y_n$。那么，模空间 $\M_{E^e/S^e}(y_0)$ 和 $\M_{E^e/S^e}(y_n)$ 的维数分别为 $n-1$ 与 $2n-1$。取其中第一个空间，考虑带参数的求值映射；这里用平行移动将 $\partial S^e$ 任意点上的 Lagrangian 边界条件与 $L_0$ 等同：
\[
\operatorname{ev}:\R\times\M_{E^e/S^e}(y_0)\longrightarrow L_0.
\]
可以证明，$\operatorname{ev}$ 在 $L_0\setminus\{y_0,\tau_V(y_0)\}$ 上的度数为奇数；缺少的这两点对应于 $z$ 趋于无穷时的极限。证明方法是先归约到闭圆盘的情形，再归约到标准局部模型；参见\bcite{131}笔记 2.17(iii)。一旦得到这一点，就容易看出 (iv) 的贡献确实是 $e\langle\cdot,\cdot\rangle$。
\SourceFigure{figure-17-10}{0.458999506}{17.10}
\end{note}
